\chapter{Introduzione all'Agentic AI nell'Information Retrieval}
\label{chap:introduzione}

\section{Contesto ed Evoluzione dell'Information Retrieval: Dal Lexical Search ai Sistemi Agentici}
\label{sec:contesto_evoluzione}

L'Information Retrieval (IR) costituisce uno dei pilastri fondamentali dell'informatica moderna, preposto all'individuazione, estrazione e indicizzazione di informazioni non strutturate da vaste collezioni documentali in risposta a uno specifico bisogno informativo espresso dall'utente (\textit{information need}). Storicamente, l'evoluzione dei paradigmi di Information Retrieval può essere scandita attraverso tre generazioni concettuali distinte:
\begin{enumerate}
    \item \textbf{Modelli Lessicali e Probabilistici}: Basati sull'ipotesi che la rilevanza di un documento sia correlata alla frequenza e alla distribuzione dei termini della query condivisi con il testo. Modelli classici quali il modello Booleano, lo spazio vettoriale (Vector Space Model con pesatura TF-IDF) e le formulazioni probabilistiche del \textit{Probabilistic Relevance Framework} culminate nell'algoritmo BM25 \cite{robertson2009probabilistic}, hanno dominato la scena per decenni grazie alla loro eccezionale scalabilità computazionale e solidità teorica. Tuttavia, tali approcci soffrono intrinsecamente del divario di vocabolario (\textit{vocabulary mismatch problem}), fallendo nell'interpretare sinonimi, relazioni tassonomiche o sfumature semantiche contestuali.
    \item \textbf{Modelli Densi e Semantici (Dense Retrieval)}: Con l'avvento dei modelli di linguaggio basati su architettura Transformer, il paradigma si è spostato verso il recupero neurale denso (Dense Passage Retrieval, Bi-Encoders e Cross-Encoders) \cite{karpukhin2020dense}. I documenti e le interrogazioni vengono proiettati in uno spazio vettoriale continuo latente $\mathbb{R}^d$, in cui la prossimità semantica è quantificata tramite misure geometriche di similarità (prodotto scalare o similarità coseno). Pur superando il problema del vocabolario, i modelli densi standard rimangono strutturalmente passivi: assumono che la risposta al bisogno informativo sia interamente contenuta all'interno di un singolo segmento testuale isolato (\textit{single-shot retrieval}) e manifestano una marcata degradazione prestazionale quando confrontati con query complesse che richiedono composizionalità logica o dati al di fuori della distribuzione di addestramento.
    \item \textbf{Retrieval-Augmented Generation (RAG) Tradizionale e i Suoi Limiti}: L'integrazione di Large Language Models (LLM) con meccanismi di recupero documentale ha dato origine ai sistemi RAG convenzionali (\textit{Naïve RAG}). In tale schema, un motore di ricerca recupera i primi $k$ passaggi rilevanti, i quali vengono concatenati all'interno del contesto dell'LLM per generare la risposta finale. Sebbene la RAG riduca le allucinazioni parametriche del modello, essa evidenzia severi limiti strutturali:
    \begin{itemize}
        \item \textit{Assenza di reattività}: Il recupero avviene in un unico passaggio cieco. Se i passaggi estratti non contengono l'evidenza risolutiva o risultano rumorosi, il generatore sintetizza risposte errate o fuorvianti senza possibilità di appello.
        \item \textit{Inadeguatezza per query multi-hop}: Quando il soddisfacimento del bisogno informativo richiede l'aggregazione di evidenze disgiunte e sequenziali (ad esempio: \textit{"Chi è il regista del film vincitore dell'Oscar 2024 e dove ha conseguito la laurea?"}), il retrieval a singolo passaggio non è in grado di formulare la query intermedia necessaria senza aver prima risolto la prima incognita.
        \item \textit{Mancanza di auto-valutazione}: Il sistema non valuta criticamente se i documenti recuperati siano effettivamente rilevanti rispetto all'intento dell'utente o se la risposta generata sia strettamente ancorata (\textit{grounded}) alle evidenze fornite.
    \end{itemize}
\end{enumerate}

Per superare queste barriere architettoniche, la ricerca più recente ha promosso la convergenza tra Information Retrieval e Agentic AI, delineando il nuovo paradigma dell'\textbf{Agentic Information Retrieval (Agentic IR)} \cite{zhang2024ai, wu2025agentic}. In questo contesto, l'Information Retrieval cessa di essere una procedura rigida e unidirezionale di classificazione passiva dei documenti e si trasforma in un processo esplorativo attivo, guidato da agenti intelligenti capaci di interagire dinamicamente con l'ambiente, orchestrare strumenti computazionali eterogenei e adattare ricorsivamente la propria strategia investigativa.

---

\section{Fondamenti e Proprietà Cardine dell'Agentic AI in Ambito IR}
\label{sec:fondamenti_agentic}

Un sistema di intelligenza artificiale viene definito \textbf{agentico} quando manifesta una capacità intrinseca di \textit{agency}, ossia l'abilità di perseguire autonomamente obiettivi complessi di alto livello mediante cicli deliberativi di percezione, pianificazione, esecuzione di azioni e riflessione critica sull'ambiente circostante.

Formalmente, un agente per l'Information Retrieval può essere modellato come un decisore operante su uno spazio di stati $\mathcal{S}$, un insieme di azioni ammissibili $\mathcal{A}$ e uno spazio di osservazioni $\mathcal{O}$. Dato un bisogno informativo espresso dall'utente mediante una query iniziale $q_0$, il ciclo operativo dell'agente si articola secondo il principio ciclico \textbf{Percezione $\rightarrow$ Ragionamento $\rightarrow$ Azione $\rightarrow$ Feedback}:
\begin{equation}
    s_t = \tau(s_{t-1}, a_{t-1}, o_t), \quad a_t \sim \pi(a \mid s_t, q_0)
\end{equation}
dove $\tau$ rappresenta la funzione di transizione dello stato interno (che memorizza la cronologia dei passaggi, le evidenze raccolte e i fallimenti precedenti), $o_t \in \mathcal{O}$ è l'osservazione restituita dall'ambiente all'istante $t$ (ad esempio, il contenuto dei passaggi estratti da una sorgente documentale), e $\pi$ è la politica deliberativa guidata dall'LLM.

Come evidenziato dalle linee guida accademiche e dalla letteratura di riferimento, un'architettura Agentic IR matura si caratterizza per quattro proprietà fondamentali:

\subsection{Autonomia Operativa e Reattività Adattiva}
A differenza dei modelli reattivi che si limitano a produrre un output token-per-token in risposta a un prompt fisso, l'agente opera con discrezionalità strategica: analizza la complessità della query, stima il proprio livello di incertezza e determina autonomamente quando il contesto informativo raccolto è sufficiente per generare la risposta finale o se siano necessarie ulteriori iterazioni investigative \cite{jeong2024adaptiverag}.

\subsection{Goal-Directed Planning e Decomposizione Gerarchica}
Di fronte a problemi informativi composti, l'agente non esegue interrogazioni casuali, ma formula piani d'azione strutturati. Attraverso paradigmi come il \textit{Least-to-Most Prompting} \cite{zhou2022leasttomost} e il framework \textit{ReAct} (Reasoning + Acting) \cite{yao2022react}, l'agente decompone una query globale complessa in una sequenza ordinata di sotto-problemi atomici (\textit{sub-queries}) dipendenti o indipendenti, risolvendo ciascun frammento mediante passaggi d'inferenza intermedi:
\begin{equation}
    q_0 \xrightarrow{\text{Decomposizione}} \{q_1, q_2, \dots, q_m\}, \quad q_i = f(q_0, o_{<i})
\end{equation}
in cui la formulazione della sotto-query $q_i$ può dipendere semanticamente dalle risposte parziali ottenute per le sotto-query precedenti $q_{<i}$.

\subsection{Integrazione e Orchestrazione di Tool Eterogenei (Tool Augmentation)}
Un agente di IR non è confinato ai pesi del proprio modello parametrico né a un singolo indice locale. Esso dispone di un registro di strumenti (\textit{Tool Registry}) che include:
\begin{itemize}
    \item Motori di ricerca lessicali su indici invertiti locali (ad es. implementazioni BM25 per documentazione verificata e ad alta velocità);
    \item API di motori di ricerca su larga scala sul Web (per il reperimento di informazioni fresche, dati di attualità o entità della coda lunga);
    \item Basi di conoscenza strutturate (Knowledge Graphs, basi di dati relazionali SQL/NoSQL).
\end{itemize}
L'agente sintetizza la chiamata allo strumento, genera gli argomenti d'interrogazione mirati, interpreta l'output restituito e incorpora i risultati nell'accumulatore di contesto.

\subsection{Autocritica e Ragionamento Auto-Riflessivo (Self-Reflection \& Grounding)}
L'agente implementa meccanismi di validazione metacognitiva per valutare la qualità dell'informazione reperita, ispirandosi ai principi di \textit{Self-RAG} \cite{asai2023selfrag}. Esso valuta sistematicamente:
\begin{enumerate}
    \item \textbf{Relevance Critique ($\text{IsRel}$)}: I passaggi recuperati contengono informazioni effettivamente pertinenti alla domanda o introducono rumore non correlato?
    \item \textbf{Supportiveness Critique ($\text{IsSup}$)}: La risposta proposta è interamente deducibile e comprovata dalle evidenze documentali estratte, oppure introduce speculazioni prive di fondamento fattuale (\textit{hallucinations})?
\end{enumerate}
Qualora i criteri di grounding non siano soddisfatti, l'agente rifiuta la risposta e re-instrada il flusso operativo verso una fase di raffinamento della query o di pianificazione correttiva.

---

\section{Definizione del Problema e Sfide Aperte nell'IR Complesso}
\label{sec:definizione_problema}

L'obiettivo cardine di questa ricerca risiede nella progettazione, implementazione e validazione empirica di un'architettura agentica autonoma per compiti di Information Retrieval complessi.

Nel recupero di informazioni su larga scala, la maggior parte delle interrogazioni del mondo reale non si presenta nella forma di semplici domande a risposta diretta su singole entità ben note (\textit{factoid single-hop queries}), bensì come interrogazioni composite che sollevano tre principali sfide aperte:

\subsection{La Sfida delle Query Multi-Hop e Comparative}
Nei benchmark avanzati di Information Retrieval e Question Answering (quali HotpotQA \cite{yang2018hotpotqa}), le domande richiedono tipicamente di individuare e concatenare informazioni distribuite su molteplici documenti non correlati a priori. Si distinguono due classi strutturali critiche:
\begin{itemize}
    \item \textbf{Bridge Questions (Domande a Ponte)}: Per identificare l'entità target finale, occorre prima determinare un'entità intermedia che funge da "ponte". Esempio: \textit{"Quale altra scuderia di Formula 1 ha guidato il pilota che ha vinto il Gran Premio di Monaco nel 1996?"}. Un motore tradizionale interrogherà l'indice con l'intera frase, trovando documenti generici sul GP del 1996 ma fallendo nel collegare l'informazione che il pilota è Olivier Panis, per poi cercare le altre scuderie guidate da Panis.
    \item \textbf{Comparison Questions (Domande Comparative Multi-Entità)}: Interrogazioni che mettono a confronto attributi di due o più entità distinte. Esempio: \textit{"Scott Derrickson ed Ed Wood erano della stessa nazionalità?"}. La query richiede di estrarre in parallelo o in sequenza la nazionalità di Derrickson (statunitense), quella di Wood (statunitense), e infine applicare un'operatore di uguaglianza booleana per concludere con successo.
\end{itemize}

\subsection{Il Divario Informativo (Information Gap) e il Semantic Drift}
Spesso gli utenti formulano interrogazioni affette da polisemia, ambiguità lessicale o sotto-specificazione. In un sistema RAG convenzionale, l'esecuzione di una ricerca su una query ambigua produce documenti fuorvianti; l'LLM, condizionato da tali documenti, genera una risposta errata che consolida il fenomeno del \textit{semantic drift} (deriva semantica). Risulta dunque indispensabile un'unità specializzata di \textit{Query Refinement} che ri-scriva, espanda termini sinonimi o disambigui le entità prima o durante il recupero.

\subsection{La Gestione del Trade-Off tra Efficienza, Allucinazioni e Copertura}
Un'ulteriore sfida critica riguarda il bilanciamento tra l'utilizzo di una Knowledge Base locale (sicura, deterministica, ad altissima velocità di consultazione ma intrinsecamente circoscritta nel vocabolario) e motori di ricerca Web aperti (vastissima copertura ma soggetti a rumore, formati eterogenei e latenza di rete). Un agente avanzato deve saper coordinare entrambe le sorgenti, sfruttando la memoria locale come primo livello di cache semantica e ricorrendo al web unicamente per colmare specifici deficit di conoscenza.

---

\section{Obiettivi del Progetto e Contributi della Ricerca}
\label{sec:obiettivi_contributi}

Il presente progetto, sviluppato nell'ambito del corso di \textit{Information Retrieval} presso l'Università degli Studi di Napoli Federico II (Laboratorio IKNOS-DIETI), affronta in modo integrato i task ufficiali delineati nella traccia d'esame (Tabella~1 del documento di specifica):
\begin{itemize}
    \item \textbf{Autonomous Query Refinement}: Analisi dell'intento dell'utente, rilevamento della complessità strutturale e decomposizione di query complesse in sotto-query atomiche (\textit{Least-to-Most decomposition}), con riscrittura e ripulitura lessicale.
    \item \textbf{Multi-Step Retrieval Planning}: Formulazione ed esecuzione di piani investigativi multi-hop mediante cicli iterativi ReAct, orchestrando passaggi di recupero progressivi e integrando evidenze parziali.
    \item \textbf{Tool-Augmented Retrieval}: Integrazione ibrida tra un indice locale probabilistico basato su BM25 (costruito su un corpus enciclopedico di riferimento) e strumenti di Web Search in tempo reale per estendere dinamicamente il perimetro di conoscenza.
    \item \textbf{Self-Reflective Validation and Fact Checking}: Verifica critica sistematica (ispirata a Self-RAG e Adaptive-RAG) dello stato di avanzamento, della pertinenza delle evidenze raccolte e della rigorosa aderenza fattuale (\textit{grounding}) della risposta proposta prima della terminazione del grafo.
\end{itemize}

\subsection{L'Architettura Proposta a Grafo di Stati (LangGraph)}
L'infrastruttura ingegnerizzata abbandona la linearità delle pipeline tradizionali in favore di un grafo orientato a stati finiti implementato mediante \textbf{LangGraph}, in cui lo stato globale dell'investigazione (\textit{Agentic State}) fluisce attraverso tre nodi agentici specializzati:
\begin{enumerate}
    \item \textbf{Validator (Adaptive Routing \& Critique Node)}: Funge da supervisore e controllore centrale del flusso. Esamina la complessità della query originale secondo i principi di \textit{Adaptive-RAG}, decide se instradarla al raffinamento o alla pianificazione diretta, e valida l'output generato mediante controlli di consistenza logica e polarità booleana.
    \item \textbf{Refiner (Autonomous Query Refinement Node)}: Specializzato nella manipolazione semantica della query. Riconosce le domande comparative e a ponte, applica la decomposizione Least-to-Most generando sotto-interrogazioni atomiche auto-consistenti e preserva la coerenza referenziale.
    \item \textbf{Planner (Multi-Step ReAct Retrieval Planning Node)}: Esegue l'esplorazione documentale vera e propria. Per ciascuna sotto-query, avvia un ciclo ReAct multi-hop alternando pensieri analitici (\textit{Thoughts}), azioni di recupero mirate (\textit{Actions}: ricerca locale BM25 o ricerca Web) e interpretazione delle evidenze restituite (\textit{Observations}), sintetizzando infine una risposta fondata sui soli fatti provati.
\end{enumerate}

\subsection{Valutazione Sperimentale Rigorosa}
Per validare empiricamente l'efficacia del framework rispetto a una baseline classica vincolata alla sola Knowledge Base locale, è stato implementato un protocollo di benchmark oggettivo sul subset di validazione di \textbf{HotpotQA} \cite{yang2018hotpotqa}, valutando le prestazioni attraverso le metriche standard della letteratura di Information Retrieval: \textit{Grounded Accuracy}, \textit{Token-level Macro F1-Score}, \textit{Exact Match (EM)}, \textit{Precision} e \textit{Recall}. I risultati sperimentali (esposti in dettaglio nel Capitolo~\ref{chap:implementazione_valutazione}) dimostrano come il coordinamento multi-agente incrementi l'accuratezza dal $26.7\%$ della baseline al $43.3\%$ del sistema integrale, confermando l'ipotesi che la combinazione di decomposizione della query, pianificazione ReAct e verifica di grounding consenta di superare drasticamente i limiti strutturali dell'IR classico.

---

\section{Struttura della Relazione}
\label{sec:struttura_relazione}

Il presente documento è organizzato in conformità con le direttive del corso di Information Retrieval e si articola nei seguenti capitoli:
\begin{itemize}
    \item \textbf{Capitolo~\ref{chap:dati_preprocessing} - Dati e Preprocessing}: Descrive i dataset di riferimento utilizzati per l'addestramento e il benchmarking (con particolare attenzione al dataset HotpotQA per compiti multi-hop e a Wikipedia come corpus documentale), illustrando le tecniche di normalizzazione testuale, indicizzazione BM25 e strutturazione dei contesti documentali.
    \item \textbf{Capitolo~\ref{chap:metodologie} - Metodologie di Retrieval e Ragionamento Agentico}: Approfondisce le fondamenta teoriche e algoritmiche adottate nel framework: il paradigma ReAct per l'esplorazione guidata da tool, l'Adaptive-RAG per la classificazione dinamica della complessità delle query, il Least-to-Most prompting per la decomposizione composizionale, e i meccanismi di Self-RAG per l'autocritica e la mitigazione delle allucinazioni.
    \item \textbf{Capitolo~\ref{chap:implementazione_valutazione} - Implementazione del Sistema, Valutazione Sperimentale e Risultati}: Espone dettagliatamente l'architettura software del framework multi-agente basato su LangGraph, la persistenza dello stato su database SQLite relazionale, la Dashboard interattiva sviluppata in Streamlit per l'ispezione real-time e l'analisi quantitativa comparativa tra il Sistema Integrale e la Baseline su HotpotQA.
    \item \textbf{Capitolo~\ref{chap:conclusioni} - Conclusioni e Sviluppi Futuri}: Riassume i principali contributi scientifici e ingegneristici del lavoro, discute i limiti intrinseci emersi dalle evidenze sperimentali e delinea le direzioni di ricerca future, incluse le possibili estensioni verso tesi di laurea magistrale in ambiti specialistici (quali l'analisi spazio-temporale o la sicurezza informatica).
\end{itemize}

